\section{Fredholm theory for a continuous positive form}
\label{sec:continuous-hermitian-fredholm}

A continuous strictly positive $(1,1)$-form still defines a uniformly
elliptic trace operator. Its range is most conveniently described in Sobolev spaces:
continuity of the coefficients suffices for the local $W^{2,p}$ estimate,
while the maximum principle identifies the kernel. The lemma is used for
continuous positive complex Hessians in
\cref{thm:continuous-hessian-adjoint-uniqueness}.

\begin{lemma}\label{lem:continuous-hermitian-fredholm}\index{Elliptic operator!continuous coefficients}
Let $\sigma$ be a continuous strictly positive real $(1,1)$-form on $X$,
and let $\mathrm dV$ be a smooth positive volume form.
For $p>2n$, the real operator
\[
 L_\sigma v=\operatorname{tr}_\sigma\ddc v,
 \qquad L_\sigma\colon W^{2,p}(X)\longrightarrow L^p(\mathrm dV),
\]
is Fredholm of index zero. Its kernel consists of the constants, and
its range has a one-dimensional annihilator in $L^{p'}(\mathrm dV)$,
where $p'=p/(p-1)$.
\end{lemma}
\begin{proof}
Fix a smooth positive Hermitian form $\eta$ and put
$L_s=(1-s)L_\eta+sL_\sigma$, $0\leq s\leq1$. On a finite holomorphic
cover these operators have uniformly positive real principal matrices
and a common modulus of continuity. They have no lower-order terms in
these coordinates. Write $d=2n$ for the real dimension.

\emph{Local estimates.} For a constant positive definite real matrix
$a_0$, the symbols $\xi_i\xi_j/(1+a_0^{kl}\xi_k\xi_l)$ satisfy the
order-zero multiplier estimates, uniformly for $a_0$ in a compact set of
positive definite matrices. Consequently
\[
 \|D^2z\|_p\leq C\bigl(\|a_0^{ij}D_{ij}z\|_p+\|z\|_p\bigr)
 \qquad(z\in C_c^\infty).
\]
Freeze the principal matrix of $L_s$ at a point and shrink the ball until
the coefficient oscillation times this constant is less than $1/2$.
The oscillation term can then be absorbed on the left. Differentiating
a cutoff adds first derivatives; the interpolation inequality
$\|Dz\|_p\leq\varepsilon\|D^2z\|_p+C_\varepsilon\|z\|_p$
absorbs them. This inequality follows by comparing $Dz$ with its
mollification at scale $\varepsilon$ and estimating the derivative of
the mollified function by $C\varepsilon^{-1}\|z\|_p$.
A partition of unity and smooth density give, uniformly in $s$,
\begin{equation}\label{eq:continuous-hermitian-wtwo}
 \|v\|_{W^{2,p}}\leq C\bigl(\|L_sv\|_p+\|v\|_p\bigr).
\end{equation}

\emph{Maximum principles.} We give the local estimate needed to pass
comparison to $W^{2,p}$. Let $L=a^{ij}D_{ij}$ on a bounded smooth
coordinate domain $\Omega$, with $a\geq\vartheta I$, and let $z$ be
smooth on $\overline\Omega$ with nonnegative boundary values and minimum
$-m<0$. Set $R=\operatorname{diam}\Omega$. Every vector $\xi$ with
$|\xi|<m/(2R)$ is the gradient at an interior lower supporting point of
$z$: minimize $z(x)-\xi\cdot(x-x_0)$, where $z(x_0)=-m$. Its boundary
values exceed $-m/2$, whereas its value at $x_0$ is $-m$. At a supporting point
$D^2z\geq0$, and the arithmetic--geometric mean inequality gives
\[
 \det D^2z\leq
 \left(\frac{\operatorname{tr}(aD^2z)}{d\vartheta}\right)^d.
\]
The area formula for the gradient, restricted to these supporting
points, therefore bounds the volume of the slope ball by
$(d\vartheta)^{-d}\int_\Omega((Lz)^+)^d$. Thus
\begin{equation}\label{eq:continuous-hermitian-local-minimum}
 -\min_\Omega z\leq C_d\vartheta^{-1}R\|(Lz)^+\|_{L^d(\Omega)}.
\end{equation}
This calculation uses no derivative of $a$. Now let $v$ be of class
$W^{2,p}$ on a neighborhood of $\overline\Omega$, and approximate it in
$W^{2,p}$ by smooth functions. Since $p>d$, the approximations converge
in $C^1(\overline\Omega)$, and their images under $L$ converge in
$L^d(\Omega)$. Adding to each approximation the negative part of its
minimum on $\partial\Omega$, which tends to zero, we may apply
\eqref{eq:continuous-hermitian-local-minimum}. Passage to the limit
shows that $Lv\leq0$ and $v|_{\partial\Omega}\geq0$ imply $v\geq0$.

The same comparison proves the strong maximum principle. If
$v\in W^{2,p}$ on a connected open set, $Lv\geq0$, and $v$ has an
interior maximum $m$, put $w=m-v\geq0$. If $w$ is positive somewhere and
zero somewhere, a small ball contained in $\{w>0\}$ and tangent to its
zero set can be chosen inside a coordinate chart. Let $c$ and $r$ be its
center and radius. On the annulus $\{r/2<|x-c|<r\}$,
\[
 \zeta(x)=\exp(-k|x-c|^2)-\exp(-kr^2)
\]
is positive and has $L\zeta>0$ for sufficiently large $k$, by uniform
ellipticity. For small $\epsilon>0$, the function $\epsilon\zeta$ is
bounded above by $w$ on both boundary spheres, hence throughout the
annulus by the comparison principle applied to $w-\epsilon\zeta$. At the tangent point
this forces a positive inward derivative of $w$, contradicting $Dw=0$ at
this minimum of the $C^1$ function $w$. Connectedness proves that $v$ is
constant.
Applied to $L_s$ on $X$, it gives $\ker L_s=\mathbb R$. If
$(L_s-1)v=0$ and $v$ has a positive maximum, apply the strong maximum
principle on the component of $\{v>0\}$ containing a maximum point,
where $L_sv=v>0$. Constancy
contradicts that equation. Applying the argument to $-v$ rules out a
negative minimum. Hence $\ker(L_s-1)=0$.

\emph{Index and range.} Compactness of $W^{2,p}\hookrightarrow L^p$
and \eqref{eq:continuous-hermitian-wtwo} now give
\begin{equation}\label{eq:continuous-hermitian-inverse}
 \|v\|_{W^{2,p}}\leq C\|(L_s-1)v\|_p
\end{equation}
uniformly in $s$. If this failed, we could normalize a contradicting sequence in
$W^{2,p}$, pass to $s_j\to s$ and a strongly $L^p$ convergent subsequence,
and apply \eqref{eq:continuous-hermitian-wtwo} to differences. Operator
norm continuity in $s$ gives strong $W^{2,p}$ convergence to a nonzero
member of $\ker(L_s-1)$, which is impossible.
The smooth operator $L_\eta-1$ has index zero and trivial kernel, hence
is an isomorphism, by smooth elliptic theory. The method of continuity
and \eqref{eq:continuous-hermitian-inverse} show that every $L_s-1$ is
an isomorphism. Finally, $L_\sigma$ differs from $L_\sigma-1$ by the
compact inclusion $W^{2,p}\hookrightarrow L^p$, so it is Fredholm of
index zero. Its one-dimensional kernel makes its range have codimension
one. $L^p$ duality gives the annihilator assertion.
\end{proof}
